\chapter{Future Work}

\section{Completeness for the probabilistic semantics}
The most immediate next step is a completeness theorem, or a proof that completeness fails for
this semantics. Either result would be valuable. A completeness proof would show that the
semantics captures exactly the inferential content of the associative Lambek calculus. A failure
of completeness would be equally informative, because it would identify the precise point at
which corpus-relative probabilistic information departs from the classical residuated behaviour.

The main obstacles are already visible. The semantics is finite, numerical, and tied to attested
fragments. Standard canonical constructions for classical Lambek models do not immediately
produce such corpus-relative valuations. Moreover, the use of maxima and minima in the
semantic clauses suggests that any completeness proof may need to build a model whose
numerical values are carefully engineered rather than merely order-theoretic.


\section{Extending the formal model}
Several extensions would make the semantics more realistic.
\begin{enumerate}[label=(\roman*)]
    \item \textbf{Out-of-vocabulary handling.} A smoothing method should assign principled values to
    unseen words or fragments.
    \item \textbf{Learning lexical probabilities.} Instead of assuming $\Lex$, one should estimate it from
    annotated or weakly supervised corpora.
    \item \textbf{Beyond observed fragments.} The present semantics could be extended from
    $\Frag(\mathcal C)$ to all of $\Sigma^+$, perhaps with backoff or compositional generalisation.
    \item \textbf{Normalisation over derivations.} A fuller probabilistic semantics could attach
    probabilities directly to derivation trees, not only to string--type pairs.
\end{enumerate}


\section{Implementation}
A concrete implementation would be a natural and worthwhile continuation of the thesis.
One plausible route is a chart parser or dynamic-programming parser for the associative Lambek
calculus with weighted lexical entries. The product clause already suggests a bottom-up dynamic
program: for each substring, compute the best score for each complex type by maximising over
binary splits. Residual values could then be computed from attested contexts or approximated by
statistics gathered during training.

A practical implementation project could proceed in four stages.
\begin{enumerate}[label=(\arabic*)]
    \item Build a lexicon of atomic type probabilities from an annotated corpus.
    \item Implement a weighted Lambek chart parser over substrings.
    \item Compare the parser's ambiguity-resolution behaviour with PCFG and probabilistic CCG
    baselines.
    \item Investigate whether the semantic values provide useful ranking information for downstream
    semantic interpretation.
\end{enumerate}

Such an implementation would also help test whether the present semantics is best interpreted as
(a) a genuine probabilistic semantics, (b) a scoring semantics for parse ranking, or (c) the first
step toward a fuller statistical Lambek model.


\section{Final remarks}
This thesis has aimed to do one thing carefully rather than many things incompletely: to turn a
promising but technically unstable sketch into a coherent formal proposal. The result is a
corpus-based probabilistic semantics for the associative Lambek calculus together with a
soundness theorem. That already provides a viable foundation for a larger research programme.
The next stage is clear: completeness, implementation, and empirical evaluation.
